%! Author = sriadityavedantam
%! Date = 4/7/26

\section{Tangent Space}

Let $X$ be an affine variety with $x \in X$.
Our goal is to define the tangent space at $x$.
We will give two definitions: the concrete and the intrinsic.

\subsection{Concrete Definition}

Assume $X \subseteq \aa^n$ and, without loss of generality, let $x = (0,\ldots,0) \in \aa^n$.
Our idea is that the tangent space at $x$, denoted $T_x$, is the union of all lines tangent to $X$ at $x$.
A tangent space exists at singular points, not just smooth points.
Different notions of tangency may not agree.
Thus calculus alone is insufficient, and we instead use intersection multiplicity.

\begin{definition}
    A line is tangent at a point $x$ if it has intersection multiplicity $\geq 2$ at $x$.
\end{definition}

\begin{proposition*}{
    Let $\ell$ be a line in $\aa^n$ passing through $x = (0,\ldots,0)$.
    Pick a point $a = (a_1,\ldots,a_n) \in \ell$.
    Then all points on $\ell$ are of the form $ta$ for $t \in \c$.
    Let $X = \{f_1 = \cdots = f_m = 0\} \subseteq \aa^n$.
    For each $i$, the restriction $f_i|_\ell$ is a polynomial in $t$.
    Since $f_i(0) = 0$, we may write
    \[
        f_i|_\ell = t^{k_i} g_i(t), \quad g_i(0) \neq 0.
    \]
}
\end{proposition*}

\begin{definition}[Multiplicity of $x$ as a root]
    The integer $k_i \in \z_{\ge 0} \cup \{\infty\}$ is called the multiplicity (order of vanishing) of $f_i|_\ell$ at $0$.
    If $f_i|_\ell \equiv 0$, then $k_i = \infty$.
\end{definition}

\begin{definition}[Intersection multiplicity]
    \[
        \mu \coloneqq \min_{1 \le i \le m} \{k_i\}.
    \]
\end{definition}

\begin{definition}[Line of tangency]
    A line $\ell$ is tangent at $x$ if $\mu \ge 2$.
\end{definition}

\begin{definition}[Tangent space]
    \[
        T_x \coloneqq \bigcup_{\ell \text{ tangent at } x} \ell.
    \]
\end{definition}

\begin{example}
    \[
        X = \{y = x^2\} \subseteq \aa^2.
    \]
    Pick $\ell$ with parameterization $(a_1 t, a_2 t)$.
    \begin{align*}
        f(x,y) &= x^2 - y, \\
        f|_\ell &= (a_1 t)^2 - a_2 t = a_1^2 t^2 - a_2 t.
    \end{align*}
    Then $\mu \ge 2 \iff a_2 = 0$.
    Thus the tangent line is $\{y = 0\}$.
\end{example}

\begin{example}
    \[
        X = \{y(y - x^2) = 0\} \subseteq \aa^2.
    \]
    Let $\ell$ be parameterized by $(a_1 t, a_2 t)$.
    \begin{align*}
        f|_\ell &= a_2 t (a_2 t - (a_1 t)^2) \\
        &= a_2^2 t^2 - a_1 a_2 t^3.
    \end{align*}
    In all cases, $\mu \ge 2$.
    Hence every line through $(0,0)$ is tangent, and
    \[
        T_{(0,0)} = \aa^2.
    \]
\end{example}

\begin{example}
    \[
        X = \{y = a x\} \subseteq \aa^2.
    \]
    \begin{align*}
        f(x,y) &= y - ax, \\
        f(a_1 t, a_2 t) &= (a_2 - a a_1)t.
    \end{align*}
    We have $\mu = \infty$ only when $a_2 = a a_1$.
    Thus the tangent line is $\{y = ax\}$.
\end{example}

\subsection{Intrinsic Definition}

Let $\o_{X,x}$ be the local ring at $x$ with maximal ideal $\m_x$.
Define
\[
    \m_x^2 \coloneqq \left\{ \sum a_i b_i : a_i, b_i \in \m_x \right\}.
\]

\begin{definition}[Cotangent space]
    The cotangent space is
    \[
        \m_x / \m_x^2.
    \]
\end{definition}

\begin{theorem*}{
    \[
        T_x \cong (\m_x / \m_x^2)^\vee.
    \]
}
\end{theorem*}

\begin{definition}[Lift]
    Given morphisms $f : X \to Y$ and $g : Z \to Y$, we say $f$ lifts through $Z$ if there exists $h : X \to Z$ such that $f = g \circ h$.
\end{definition}

\begin{remark}
    The cotangent space is a vector space over the residue field $\o_{X,x}/\m_x \cong \c$, and its dual is the tangent space.
\end{remark}

\begin{proposition*}{
    Let $X \subseteq \aa^n$ be affine with defining equations $f_1,\ldots,f_m$.
    Then
    \[
        T_x = \{ a \in \aa^n : d_x f_i(a) = 0,\ \forall i \}.
    \]
}
\end{proposition*}

\begin{example}
    If $X = \aa^1$, then $\c[\aa^1] = \c[t]$, $\m_0 = (t)$, and $\m_0^2 = (t^2)$.
    Thus
    \[
        \m_0 / \m_0^2 \cong \c,
    \]
    so $T_0 \aa^1 \cong \c$.
\end{example}